% =====================================================================
\section{\LaTeX{} und KI}
% =====================================================================

\begin{frame}{Wofür KI bei \LaTeX{} hilfreich ist}
  \begin{tcbraster}[raster columns=4, raster equal height, raster column skip=8pt]
    \begin{karte}{Fehler}{verstehen}
      Log-Ausschnitt einfügen und erklären lassen, was \LaTeX{} meint.
    \end{karte}
    \begin{karte}{Bilder}{zu Formeln}
      Foto einer Rechnung oder Formel in \LaTeX-Code umwandeln.
    \end{karte}
    \begin{karte}{Gerüste}{erzeugen}
      Tabellen, Präambeln, Plot-Grundgerüste, einfache TikZ-Grafiken.
    \end{karte}
    \begin{karte}{Befehle}{nachschlagen}
      \glqq Wie setze ich \dots?\grqq{} statt lange Doku zu durchsuchen.
    \end{karte}
  \end{tcbraster}
  \vspace{0.8em}
  \begin{merke}
    KI ersetzt nicht das Verständnis der Grundlagen: Wer weiß, was \texttt{align} oder \texttt{\textbackslash label} tun, erkennt auch, wann die Antwort falsch ist.
  \end{merke}
\end{frame}

\begin{frame}{KI direkt in Overleaf}
  \begin{columns}[T]
    \begin{column}{0.53\textwidth}
      \begin{itemize}
        \setlength{\itemsep}{0.35em}
        \item \textbf{Error Assist:} erklärt Fehler und schlägt Lösungen vor
        \item \textbf{Equation Generator:} Formel in Worten beschreiben, \LaTeX{} bekommen
        \item \textbf{Table Generator:} Tabellen aus einer Beschreibung
        \item \textbf{TeXGPT:} \LaTeX-Code aus normaler Sprache im Editor
        \item \textbf{Writefull:} Sprachvorschläge für wissenschaftliches Englisch
      \end{itemize}
    \end{column}
    \begin{column}{0.41\textwidth}
      \begin{kasten}{Kostenloser Plan}
        Begrenztes tägliches Kontingent für die KI-Funktionen. Der KI-Chat-Assistent ist nicht enthalten.
      \end{kasten}
      \vspace{0.4em}
      {\scriptsize\color{grau}Stand: Oktober 2026, Details unter \texttt{docs.overleaf.com}}
    \end{column}
  \end{columns}
\end{frame}

\begin{frame}{Grenzen und Fallstricke}
  \begin{tcbraster}[raster columns=2, raster equal height, raster column skip=12pt, raster row skip=8pt]
    \begin{karte}[akzent!70!red]{Erfundene Befehle}{Halluzinationen}
      Pakete und Befehle, die es nicht gibt. Immer kompilieren und prüfen.
    \end{karte}
    \begin{karte}[akzent!70!red]{Erfundene Literatur}{BibTeX}
      Quellen können frei erfunden sein. Jeden Eintrag selbst nachprüfen.
    \end{karte}
    \begin{karte}[grau]{Komplexe Grafiken}{TikZ}
      Koordinaten und Layout gehen oft schief, Nacharbeit ist die Regel.
    \end{karte}
    \begin{karte}[grau]{Übungsblätter}{Eigenleistung}
      Hilfe beim Setzen ist kein Problem, die Physik sollte von euch sein.
    \end{karte}
  \end{tcbraster}
\end{frame}

\begin{frame}{Live-Demo}
  \zeittag{ca.\ 2 Min.}
  \begin{enumerate}
    \setlength{\itemsep}{0.6em}
    \item Handschriftliche Formel fotografieren
    \item Von einer KI in \LaTeX-Code umwandeln lassen
    \item In Overleaf einfügen und kompilieren
    \item Absichtlich einen Fehler einbauen und mit Error Assist erklären lassen
  \end{enumerate}
  \vspace{0.8em}
  \begin{merke}
    Ergebnis immer mit dem Original vergleichen: Indizes, Vorzeichen und Faktoren werden gerne verwechselt.
  \end{merke}
\end{frame}
